% Template for post or presentation abstracts submitted to ILAS 2022
% 1. Fill in the presenter's information (name, affiliation, email
% address) and talk information (title of presentation, abstract).
% 2. Compile to make sure there are no errors.
% 3. Save the .tex file for your abstract with a distinctive name,
% ideally "FamilynameGivenname.tex".
% 4. Upload your .tex file to https://clr.ie/129557
%       
% * Please keep your LaTeX commands simple, and avoid the use of
%    user-defined macros.
% * Include details of co-authors and funding acknowledgements after the abstract.
% * Upload only the LaTeX file (with .tex extension).
% * Questions: Contact Niall Madden (Niall.Madden@NUIGalway.ie)

\documentclass[a4paper]{amsart} 
\usepackage{amssymb} 
\usepackage{hyperref}
\pagestyle{empty}
\date{} 

%% IGNORE THE NEXT 4 LINES
\makeatletter % 
\def\labelenumi{\theenumi.}
\def\theenumi{\@arabic\c@enumi}
\makeatother
%% STOP IGNORING NOW

\begin{document}

\title{Companion pencils for scalar (and matrix) polynomials in the monomial basis}
\author{Fernando De Ter\'an} %Presenting author only
\address{Universidad Carlos III de Madrid} % Affiliation of presenting author
\email{fteran@math.uc3m.es} % Email address of presenting author only

\maketitle

% The abstract  ***no more than one page***

In this talk, we consider general companion pencils for scalar polynomials (given in the monomial basis) over arbitrary fields. More precisely, if 
$$
p(z)=\sum_{i=0}^n a_iz^{i}
$$ 
is a scalar polynomial, with $a_i\in\mathbb F$, for $0\leq i\leq n$, and $\mathbb F$ being an arbitrary field, then a {\em companion pencil} is of the form $L(z)=A+z B$, with $A$ and $B$ being $n\times n$ matrices with entries in the ring  of polynomials in $a_0,\hdots,a_n$ (namely $\mathbb F[a_0,\hdots,a_n]$) satisfying
$$
\det L(z)=\alpha p(z),\qquad\mbox{for some $\alpha\in\mathbb F$}.
$$
We will first show several well-know classes of companion pencils, and then we will present some theoretical results about general companion pencils, like:
\begin{itemize}
\item The Smith form of every companion pencil is $\left[\begin{smallmatrix}I_{n-1}&0\\0&\frac{1}{a_n}p(z)\end{smallmatrix}\right]$.
\item Companion pencils are nonderogatory.
\end{itemize}

We will also pay attention to the sparsity. In particular, by imposing some natural restrictions on the entries, we determine
the smallest possible number of nonzero entries in any companion pencil.


If time permits, we will also show how the notion of companion pencil is extended to matrix polynomials, and analyze some of the previous questions for this notion as well.

Most of this talk is based on \cite{dh-cortona} and \cite{dh-racsam}.


\bigskip

%% Coauthor details here (optional - delete if not applicable)
% and funding acknowledgment (optional - delete if not applicable)
\emph{This is joint work with {\bf Carla Hernando}.
Supported by Ministerio de Econom\'ia y Competitividad of Spain through grants MTM2017-
90682-REDT and MTM2015-65798-P, and by Agencia Estatal de Investigaci\'on of Spain through grant PID2019-106362GB-I00/AEI/10.13039/501100011033.}

\begin{thebibliography}{1}
\bibitem{dh-cortona} 
F. De Ter\'an, C. Hernando.
\newblock A class of quasi-sparse companion pencils.
\newblock In: {\em Structured Matrices in Numerical Linear Algebra: Analysis, Algorithms and Applications}, Bini, D.A., Di Benedetto, F., Tyrtyshnikov, E., Van Barel, M. (Eds.). INdAM series, Springer (2018)157-179.
 

\bibitem{dh-racsam}
F. De Ter\'an, C. Hernando.
\newblock A note on generalized companion pencils.
\newblock {\em  RACSAM}, 114:8  (2020).

\end{thebibliography}


\end{document}


% Please leave the next bit alone  
%%% Local Variables: 
%%% mode: latex
%%% TeX-master: "../ILAS-2022"
%%% End:
